
% Theorem System
% The following boxes are provided:
%   Definition:     \defn 
%   Theorem:        \thm 
%   Lemma:          \lem
%   Corollary:      \cor
%   Proposition:    \prop   
%   Claim:          \clm
%   Fact:           \fact
%   Proof:          \pf
%   Example:        \ex
%   Remark:         \rmk (sentence), \rmkb (block)
% Suffix
%   r:              Allow Theorem/Definition to be referenced, e.g. thmr
%   p:              Add a short proof block for Lemma, Corollary, Proposition or Claim, e.g. lemp
%                   For theorems, use \pf for proof blocks

% ======= Real examples : 

% \defn{Definition Name}{
%     A defintion.
% }

% \thmr{Theorem Name}{mybigthm}{
%     A theorem.
% }

% \lem{Lemma Name}{
%     A lemma.
% }

% \fact{
%     A fact.
% }

% \cor{
%     A corollary.
% }

% \prop{
%     A proposition.
% }

% \clmp{}{
%     A claim.
% }{
%     A reference to Theorem~\ref{thm:mybigthm}
% }

% \pf{
%    proof.

% \ex{
%     some examples. 
% }

% \rmk{
%     Some remark.
% }

% \rmkb{
%     Some more remark.
% }
\usepackage{xcolor}

% Define custom colors
\definecolor{defscol}{HTML}{ecd8d7} %For definitions
\definecolor{asumscol}{HTML}{ecd8d7} %For Assumptions

\definecolor{rmkscol}{HTML}{313160} %For remarks
\definecolor{exmscol}{HTML}{e04b52} %For examples

\definecolor{lemscol}{HTML}{2c3943} %For Lemmes
\definecolor{thmscol}{RGB}{230,120,20}
\definecolor{prpscol}{HTML}{8A9A5B}
\definecolor{corscol}{HTML}{dfd9fd} %For corrolaries

\definecolor{clmscol}{HTML}{165c58} %For claims
\definecolor{facscol}{HTML}{28a8a1} %For facts



% ============================
% Definition
% ============================
\newtcbtheorem[number within=section]{mydefinition}{Definition}
{
    enhanced,
    frame hidden,
    titlerule=0mm,
    toptitle=1mm,
    bottomtitle=1mm,
    fonttitle=\sffamily\bfseries,
    coltitle=black,
    colbacktitle=defscol!50!white,
    colback=defscol!20!white,
}{defn}

\NewDocumentCommand{\defn}{m+m}{
    \begin{mydefinition}{#1}{}
        #2
    \end{mydefinition}
}

\NewDocumentCommand{\defnr}{mm+m}{
    \begin{mydefinition}{#1}{#2}
        #3
    \end{mydefinition}
}

% ============================
% Problem Setup
% ============================
\newtcbtheorem[number within=section]{myproblem}{Problem Setup}
{
    enhanced,
    frame hidden,
    titlerule=0mm,
    toptitle=1mm,
    bottomtitle=1mm,
    fonttitle=\sffamily\bfseries,
    coltitle=black,
    colbacktitle=defscol!50!white,
    colback=defscol!20!white,
}{defn}

\NewDocumentCommand{\probs}{m+m}{
    \begin{myproblem}{#1}{}
        #2
    \end{myproblem}
}

\NewDocumentCommand{\problems}{mm+m}{
    \begin{myproblem}{#1}{#2}
        #3
    \end{myproblem}
}

% ============================
% Message
% ============================

% Define a counter that resets at each \section
\newcounter{msgcounter}[section]
\renewcommand{\themsgcounter}{\thesection.\arabic{msgcounter}}

% Define your custom message box (no need to change this part)
\newtcolorbox{custommessagebox}[2][]{
    enhanced,
    coltitle=black,
    colbacktitle=defscol!50!white,
    colback=defscol!20!white,
    fonttitle=\sffamily\bfseries,
    title=#1,        % title passed in
    toptitle=1mm,
    bottomtitle=1mm,
    titlerule=0mm,
    frame hidden,
    #2
}

% Define the \msg command with auto-numbered title based on section
\NewDocumentCommand{\msg}{mm+m}{
    \refstepcounter{msgcounter} % increments the counter and allows referencing
    \begin{custommessagebox}[#1~\themsgcounter:~\textbfs{#2}]{}
        #3
    \end{custommessagebox}
}





% ============================
% Research Problem
% ============================
\newtcbtheorem{myresearchproblem}{Research Problem}
{
    enhanced,
    frame hidden,
    titlerule=0mm,
    toptitle=1mm,
    bottomtitle=1mm,
    fonttitle=\sffamily\bfseries,
    coltitle=black,
    colbacktitle=defscol!40!white,
    colback=defscol!20!white,
}{defn}

\NewDocumentCommand{\probr}{m+m}{
    \begin{myresearchproblem}{#1}{}
        #2
    \end{myresearchproblem}
}

\NewDocumentCommand{\problemr}{mm+m}{
    \begin{myresearchproblem}{#1}{#2}
        #3
    \end{myresearchproblem}
}


% ============================
% Assumption
% ============================
\newtcbtheorem[use counter from=mydefinition]{myassumption}{Assumption}
{
    enhanced,
    frame hidden,
    titlerule=0mm,
    toptitle=1mm,
    bottomtitle=1mm,
    fonttitle=\sffamily\bfseries,
    coltitle=black,
    colbacktitle=asumscol!40!white,
    colback=asumscol!20!white,
}{asum}

\NewDocumentCommand{\asum}{m+m}{
    \begin{myassumption}{#1}{}
        #2
    \end{myassumption}
}

\NewDocumentCommand{\asumr}{mm+m}{
    \begin{myassumption}{#1}{#2}
        #3
    \end{myassumption}
}

% ============================
% Theorem
% ============================

% Custom theorem counter that resets every chapter (not section!)
% Number theorems within sections
\numberwithin{equation}{section}

% Set up a custom counter for theorems
\newcounter{mytheorem}[section]
\renewcommand{\themytheorem}{\Roman{chapter}.\thesection.\arabic{mytheorem}}

% Define the theorem environment with custom label
\newtcbtheorem[use counter=mytheorem, number within=section]{mytheorem}{Theorem}{
    enhanced,
    frame hidden,
    titlerule=0.0mm,
    fonttitle=\sffamily\bfseries,
    coltitle=black,
    colbacktitle=thmscol!20!white,
    colback=thmscol!8!white,
}{thm}

% Proof environment (page-break friendly)
\makeatletter
\newenvironment{thmpf}{%
  \par\noindent{\it \textbfs{Proof for Theorem.}}\par
  % Key bits: breakable + enhanced jigsaw + allowdisplaybreaks
  \tcolorbox[
    blanker,
    breakable,
    enhanced jigsaw,
    parbox=false,
    colback=white,
    left=5mm,
    before upper={\parindent15pt\relax\allowdisplaybreaks},
    after skip=10pt,
    borderline west={1mm}{0pt}{thmscol!20!white}
  ]%
}{%
  \textcolor{thmscol!18!white}{\hbox{}\nobreak\hfill$\blacksquare$}%
  \endtcolorbox
}
\makeatother



% Combined theorem macro
\newcommand{\thm}[3]{%
    \begin{mytheorem}{#1}{#2}%
        #3%
    \end{mytheorem}%
     \par
}


% Combined theorem+proof macro
\newcommand{\thmp}[4]{%
    \begin{mytheorem}{#1}{#2}%
        #3%
    \end{mytheorem}%
     \par
    \begin{thmpf}%
        #4%
    \end{thmpf}%
    
}


% ============================
% Lemma (breakable-ready)
% ============================
\newtcbtheorem[
  use counter=mytheorem,
  number within=section
]{mylemma}{Lemma}{
  enhanced,
  % Add 'breakable' here ONLY if you want the *lemma box* to split across pages:
  % breakable,
  frame hidden,
  titlerule=0mm,
  toptitle=1mm,
  bottomtitle=1mm,
  fonttitle=\sffamily\bfseries,
  coltitle=black,
  colbacktitle=corscol!50!white,
  colback=corscol!20!white,
}{lem}

% \lem{title}{label}{content}
%             ^^^^^  ^^^^^^^
% Make 'content' long: +m
\NewDocumentCommand{\lem}{m m +m}{%
  \begin{mylemma}[label=#2]{#1}{}%
    #3%
  \end{mylemma}%
}

% Proof environment for lemma (kept breakable)
\newenvironment{lempf}{%
  \par\noindent{\it \textbfs{Proof for Lemma.}}\par
  \tcolorbox[
    blanker,breakable,left=5mm,parbox=false,
    before upper={\parindent15pt},
    after skip=10pt,
    colback=white,
    borderline west={1mm}{0pt}{corscol!50!white}
  ]%
}{%
  \textcolor{corscol!50!white}{\hbox{}\nobreak\hfill$\blacksquare$}%
  \endtcolorbox
}

% \lemp{title}{label}{content}{proof}
%                             ^^^^^  make both body args long: +m +m
\NewDocumentCommand{\lemp}{m m +m +m}{%
  \begin{mylemma}[label=#2]{#1}{}%
    #3%
  \end{mylemma}%
  \par
  \begin{lempf}%
    #4%
  \end{lempf}%
}

% ============================
% Corollary
% ============================

\newtcbtheorem[use counter from=mydefinition]{mycorollary}{Corollary}
{
    enhanced,
    frame hidden,
    titlerule=0mm,
    toptitle=1mm,
    bottomtitle=1mm,
    fonttitle=\sffamily\bfseries,
    coltitle=black,
    colbacktitle=corscol!40!white,
    colback=corscol!20!white,
}{cor}

\NewDocumentCommand{\cor}{+m}{
    \begin{mycorollary}{}{}
        #1
    \end{mycorollary}
}

% Proof environment — page-break friendly
\makeatletter
\newenvironment{corpf}{%
  \par\noindent{\it \textbfs{Proof for Corollary.}}\par
  \tcolorbox[
    blanker,
    breakable,          % allow the box to split across pages
    enhanced jigsaw,    % robust page breaking (esp. with math)
    colback=white,
    parbox=false,
    left=5mm,
    before upper={\parindent15pt\allowdisplaybreaks}, % let long displays break
    after skip=10pt,
    borderline west={1mm}{0pt}{corscol!40!white}
  ]%
}{%
  \textcolor{corscol!40!white}{\hbox{}\nobreak\hfill$\blacksquare$}%
  \endtcolorbox
}
\makeatother

\NewDocumentCommand{\corp}{m+m+m}{
    \begin{mycorollary}{{{#1}}}{}
        #2
    \end{mycorollary}

    \begin{corpf}
     {   #3}
    \end{corpf}
}


% Named proposition without proof
\NewDocumentCommand{\cornp}{m+m}{ 
    \begin{mycorollary}{{{#1}}}{}
        #2
    \end{mycorollary}
}



% ============================ % 
% Proposition % 
% ============================
% Define proposition environment with custom label and style
\newtcbtheorem[use counter=mytheorem, number within=section]{myproposition}{Proposition}
{ 
  enhanced,
  frame hidden,
  titlerule=0mm,
  toptitle=1mm,
  bottomtitle=1mm,
  fonttitle=\sffamily\bfseries,
  coltitle=black,
  colbacktitle=prpscol!18!white,
  colback=prpscol!8!white,
}
{prop}

% Proposition with label and no proof
\NewDocumentCommand{\proppnp}{m m m}{%
  \begin{myproposition}[label=#2]{#1}{}%
    #3%
  \end{myproposition}%
}

% Proof environment for propositions
\newenvironment{proppf}
{ 
  {\noindent{\it \textbfs{Proof for Proposition.}}} 
  \tcolorbox[blanker,breakable,left=5mm,parbox=false, colback=white,
    before upper={\parindent15pt}, after skip=10pt, 
    borderline west={1mm}{0pt}{prpscol!18!white}] 
}
{ 
  \textcolor{prpscol!18!white}{\hbox{}\nobreak\hfill$\blacksquare$} 
  \endtcolorbox 
}

% Proposition with label and proof
\NewDocumentCommand{\proppp}{m m m m}{%
  \begin{myproposition}[label=#2]{#1}{}%
    #3%
  \end{myproposition}%
  \begin{proppf}%
    #4%
  \end{proppf}%
}

% Add this after loading the cleveref package



% ============================
% Claim
% ============================

\newtcbtheorem[use counter from=mydefinition]{myclaim}{Claim}
{
    enhanced,
    frame hidden,
    titlerule=0mm,
    toptitle=1mm,
    bottomtitle=1mm,
    fonttitle=\sffamily\bfseries,
    coltitle=black,
    colbacktitle=clmscol!40!white,
    colback=clmscol!20!white,
}{clm}

\NewDocumentCommand{\clm}{m+m}{
    \begin{myclaim*}{#1}{}
        #2
    \end{myclaim*}
}

\newenvironment{clmpf}{
	{\noindent{\it \textbfs{Proof for Claim.}}}
	\tcolorbox[blanker,breakable,left=5mm,parbox=false,
    before upper={\parindent15pt},
    after skip=10pt,
	borderline west={1mm}{0pt}{clmscol!40!white}]
}{
    \textcolor{clmscol!40!white}{\hbox{}\nobreak\hfill$\blacksquare$} 
    \endtcolorbox
}

\NewDocumentCommand{\clmp}{m+m+m}{
    \begin{myclaim*}{#1}{}
        #2
    \end{myclaim*}

    \begin{clmpf}
        #3
    \end{clmpf}
}

% ============================
% Fact
% ============================

\newtcbtheorem[use counter from=mydefinition]{myfact}{Fact}
{
    enhanced,
    frame hidden,
    titlerule=0mm,
    toptitle=1mm,
    bottomtitle=1mm,
    fonttitle=\sffamily\bfseries,
    coltitle=black,
    colbacktitle=facscol!40!white,
    colback=facscol!20!white,
}{fact}

\NewDocumentCommand{\fact}{+m}{
    \begin{myfact}{}{}
        #1
    \end{myfact}
}

% ============================
% Proof
% ============================

\NewDocumentCommand{\pf}{+m}{
    \begin{proof}
        [\noindent\textbfs{Proof.}]
        #1
    \end{proof}
}

% ============================
% Example
% ============================


\newenvironment{myexample}{
    \tcolorbox[blanker,breakable,left=5mm,parbox=false,
    colback=white,
    before upper={\parindent15pt},
    after skip=10pt,
	borderline west={1mm}{0pt}{clmscol!25!white}]
}{
    \textcolor{clmscol!25!white}{\hbox{}\nobreak\hfill$\blacksquare$} 
    \endtcolorbox
}

\NewDocumentCommand{\exm}{m+m}{\vspace{0.5cm}
{\noindent{\textbfs{Example: #1 }}}
    \begin{myexample}
	
        #2
    \end{myexample}
}


% ============================
% Remark
% ============================


\newenvironment{myremark}{
    \tcolorbox[blanker,breakable,left=5mm,parbox=false,
    colback=white,
    before upper={\parindent15pt},
    after skip=10pt,
	borderline west={1mm}{0pt}{rmkscol!20!white}]
}{
    \textcolor{rmkscol!20!white}{\hbox{}} 
    \endtcolorbox
}

\NewDocumentCommand{\rmkb}{+m}{\vspace{0.5cm}{\par\noindent{\textbfs{Remark.}}}
    \begin{myremark}
        #1
    \end{myremark}
}



\tcbuselibrary{theorems,skins,hooks}
\newtcbtheorem[number within=section]{Theorem}{Theorem}
{%
	enhanced,
	breakable,
	colback = mytheorembg,
	frame hidden,
	boxrule = 0sp,
	borderline west = {2pt}{0pt}{mytheoremfr},
	sharp corners,
	detach title,
	before upper = \tcbtitle\par\smallskip,
	coltitle = mytheoremfr,
	fonttitle = \sffamily\bfseries\sffamily,
	description font = \mdseries,
	separator sign none,
	segmentation style={solid, mytheoremfr},
}
{th}










% % Old styles hh 

% %--------------------------------------------------------------------------
% % 		THEOREMES STYLE
% %--------------------------------------------------------------------------

% %-------		DEFINITION 		-------	
% \newcounter{defo}[chapter]
% \newenvironment{defi}[1]{\refstepcounter{defo} 
% \begin{tcolorbox}[colback=yellow!20!white,colframe=yellow!15!black,title= \textbfs{Définition \thechapter \ $\blacklozenge$ \thedefo \ | #1}]}{\end{tcolorbox}}

% %-------		THEOREME 		-------	
% \newcounter{th}[chapter]
% \newenvironment{thm}[1]{\refstepcounter{th}
% \begin{tcolorbox}[colback=mycolor!10,colframe=mycolor!10!black!80,title=\textbfs{ Théorème \thechapter \ $\blacklozenge$ \theth \ | #1}]}{\end{tcolorbox}}

% %-------		PROPOSITION 		-------	
% \newcounter{prop}[chapter]
% \newenvironment{propt}[1]{\refstepcounter{prop}
% \begin{tcolorbox}[colback=mycolor!5,colframe=mycolor!10!linkscolor!80 ,title=\textbfs{ Proposition \thechapter \ $\blacklozenge$ \theprop \ | #1}]}{\end{tcolorbox}}

% %-------		COROLLAIRE		-------
% \newcounter{cor}[chapter]
% \newenvironment{corr}[1]{\refstepcounter{cor}
% \begin{tcolorbox}[colback=mycolor!2,colframe=mycolor!10!linkscolor!40 ,title= Corolaire \thechapter \ $\blacklozenge$ \thecor \ | #1]}{\end{tcolorbox}}

% %-------		LEMME			-------
% \newcounter{lem}[chapter]
% \newenvironment{lemme}[1]{\refstepcounter{lem}
% \begin{tcolorbox}[colback=mycolor!10!blue!2,colframe=mycolor!50!blue!30,title=\textbfs{Lemme \thechapter \ $\blacklozenge$ \thelem \ | #1}]}{\end{tcolorbox}}

% %-------		METHODE 			-------
% \newcounter{met}[chapter]
% \newenvironment{meth}[1]{\refstepcounter{met}
% \begin{tcolorbox}
% [enhanced jigsaw,breakable,pad at break*=1mm,
%  colback=red!20!white,boxrule=0pt,frame hidden,
%  borderline west={1.5mm}{-2mm}{red}] \color{red}
% {\textbfs{Méthode \thechapter \ $\blacklozenge$ \themet \ | #1} } \color{black} \\ } {\end{tcolorbox}}

% %-------		REMARQUE 			-------
% \newcommand{\NB}[1]{
% \ \\
% \begin{tabular}{p{0.05\textwidth}p{0.80\textwidth}}
% \hline
% \vspace{-0.1cm} \includegraphics[scale=0.03]{./system/IDEA.png} & 	#1\\
% \hline
% \end{tabular}
% \ 
% \newline \ \newline
%  }

% %-------		EXEMPLE  		-------
% \newenvironment{exm}{ \begin{tcolorbox}
% [enhanced jigsaw,breakable,pad at break*=1mm,
%  colback=cyan!20!white,boxrule=0pt,frame hidden,
%  borderline west={1.5mm}{-2mm}{cyan}] \color{cyan}
% {\textbfs{Exemple} } \color{black} \\ } {\end{tcolorbox}}



 % ============================
% Unnumbered Message
% ============================

\NewDocumentCommand{\msgu}{m m +m}{
    \begin{custommessagebox}[
        \IfBlankTF{#2}
            {#1}
            {#1:~\textbfs{#2}}
    ]{}
        #3
    \end{custommessagebox}
}